\section{Comparison and boundary geometry}
\subsection{Adaptive discrimination and description}
Salek--Hayashi--Winter \cite[Sections~II--V]{SHW68} study asymptotic
binary channel-discrimination exponents. Their results include equal
adaptive and nonadaptive exponents for classical--quantum channels,
and separations for some entanglement-breaking channels. Our object
is instead a covering number of a supplied family in the finite-use
metric $d_N$. Its lower packing may use a different binary test for
each pair; it does not supply a universal measurement which identifies
every codeword. The exact reduction in
Lemma~\ref{lem:productreplacer68} uses input erasure, not merely
entanglement breaking, and implies no general absence of an adaptive
advantage for entanglement-breaking channels.

Cooney--Mosonyi--Wilde \cite{CMW70} is a direct antecedent for the
input-erasure mechanism. Its two-replacer discussion and asymptotic
channel-versus-replacer exponents are distinct from an exact finite-use
covering problem. The retraction of Theorem~\ref{thm:retraction70} concerns
admissible centres, not an improvement to those discrimination exponents.
The coherent amplification used here likewise has the established unitary
discrimination antecedents \cite{Acin70,DFY70}.

The two-point programme method has a metrological antecedent in
\cite{DKG67}. Likewise the half-parameter scale has antecedents in
quantum population compression \cite{YBCH68}; there the input consists
of unknown physical copies and the memory interface may include
quantum storage. Here the encoder is given classical matrix data and
the decoder specifies an operation or state. These are different
coding tasks, even when a logarithmic coefficient has the same form.

Mele--Bittel \cite{MB67} and Oufkir--Girardi \cite{OG67} investigate
learning unknown channels in diamond distance. The present encoder
receives its target and asserts no query or sample bound. The
reference-sensitive physical classical-simulation problem of
Naik--Zartab--Gisin--Banik \cite{NZGB67} is different again: a rational
description with a diamond guarantee remains a description of a
quantum operation, not a finite-classical-message implementation
accepting an unknown entangled input. No bound here is asserted for
the hardware dimension of a universal programmable processor.

\subsection{Ranks, tangent directions and remaining joint regimes}
For a general instrument $J$, let $P_y$ project onto the support of
$J_y$. The one-sided tangent cone of the legal instrument body is
\[
 T_{\mathcal C}(J)=\overline{\{t(L-J):t\ge0,\ L\in\mathcal C\}}.
\]
Every element $H$ of this cone satisfies the marginal constraint and
$P_y^\perp H_yP_y^\perp\ge0$. These necessary conditions identify
where positivity becomes active; by themselves they assert no reuse
modulus. A boundary classification must examine the metric of
repeated use along support-changing directions and along reversible
coherent directions, rather than count affine coordinates alone.
Proposition~\ref{prop:boundaryphase67} shows that a coherent unitary
phase can have scale $N^{-1}$, whereas the preparation family in
Theorem~\ref{thm:rankentropy68} has an $N^{-1/2}$ coding scale even
at rank-deficient outputs. In that preparation family the Choi ranks
are $d\,\operatorname{rank}\sigma_y$, and the rank bounds in the
theorem are output ranks. They are not arbitrary Choi-rank strata of
all instruments.

Theorem~\ref{thm:coherententropy70} now gives a joint law in the
non-input-erasing tensor family $UXU^*\otimes\sigma_y$. It includes a
full projective unitary manifold and support-changing preparation blocks,
but no input-dependent measurement probabilities. Its mixed coefficient
$u+v/2$ is established by two explicit adaptive moduli and matching product
packings, not by a tangent-dimension count alone. These qualifications are
essential: it is not a classification of all coherent/noisy couplings.

Theorems~\ref{thm:jointcode68} and~\ref{thm:rankentropy68} settle the
joint $N,\delta$ dependence in their specified fixed-dimensional
families. They do not assert dimension-uniform constants, a uniform
limit as the positive margin tends to zero, or a full tangent-cone
classification for general disturbing instruments. The older
exponential-accuracy width result retains its multiplicative gap in
the separate label-width resource. None of these description laws
is a mutable-workspace converse for general variable-input
trajectory simulation.


\subsection{Input dependence and a common full-range witness}
The fixed-readout family of Theorem~\ref{thm:cqentropy71} is a
classical--quantum channel with a fixed quantum input measurement.
Its $N$-slot programme selects conditional row states and therefore
handles input-dependent probabilities. This is within the channel
families considered by Salek--Hayashi--Winter
\cite[Sections~II and V]{SHW68} for asymptotic binary discrimination.
Their exponent comparisons are direct antecedents, but are not a
finite-use family covering number or a rational rank-aware code.
We do not infer equality between the row-product upper bound and the
adaptive distance, nor a finite-use prohibition on adaptive advantage.
The outcome/state rows may be noncommuting and have zero blocks.
The measured basis remains fixed, so a varying measurement basis and
general coherent/noisy couplings are not parameterized by $V$.

The common-estimator argument in Section~\ref{sec:fullerror71} uses
standard informationally complete tomography \cite{PLS66} and an
elementary event-counting inequality. Its point is to control covering
multiplicity at error greater than one, where the earlier requirement
of pairwise distance greater than $2\delta$ is unavailable. This new
witness is deliberately common to all packing points; the local
small-error proof still uses pair-dependent binary tests. Neither
witness turns the supplied-description code into an unknown-instrument
learning result. The full-range extension concerns half-dimensional
coding scales and does not extend the coherent $N^{-1}$ scale by fiat.
The author-side comparisons here identify premises and conclusions;
they do not constitute an independent exhaustive priority review.


\subsection{Environment-state antecedents and the observable extension}
Das--Wilde \cite[Definition~3 and Proposition~2]{DW72} formulate a
common interaction for an environment-parametrized memory cell and
reduce adaptive reading to processing its environment states. This is
a direct antecedent of our all-row programme upper bound. It does not
by itself provide one-use seizing of all those row states.
Wilde--Berta--Hirche--Kaur \cite[Definition~36 and Theorem~37]{WBHK72}
give finite-use discrimination equalities for environment-seizable
channels. Wang--Wilde \cite[Section~III.A]{WW72} formulate their
channel-box equivalence under common transformations. We use these
state-processing principles explicitly, and do not claim to introduce
either environment parametrization or seizing. Theorem~\ref{thm:seizablecover72}
records a family-uniform ambient left inverse and transfers arbitrary
covering centres; the rank-aware rational code then supplies a
finite-use family covering law. The Pauli/Bell protocol is a standard
example of the premise, not a new source of discrimination advantage.

The imported theorem \cite[Corollary~1 and Theorem~2]{PPKKO72}
determines the adaptive multiple-query distance of ordered projective
measurements. It supplies the linear-$N$ readout modulus in
Theorem~\ref{thm:readoutentropy72}. Our contribution in this use is the
mixed family covering, direct rational projection chart and conditional
rank-preserving code, not the measurement-discrimination formula.
Its flag dimension $d^2-d$ differs from the projective-unitary dimension
$d^2-1$ in the retained coherent tensor theorem.

Varying readout is now treated when the readout remains observable,
and uniformly seizable families have an exact covering reduction.
A further boundary problem is to classify nonrecoverable varying
readout with overlapping conditional outputs, or coupled reversible
and support-changing directions lacking a common ambient seizing map.
The exact equal-row collapse in Remark~\ref{rem:unflagged72} shows
that this requires an identifiability analysis before assigning any
basis exponent. The present two theorems and the earlier full-error
half-dimensional laws do not assert a general tangent-cone classification
or uniform constants in growing dimension.


\subsection{Finite-distance noise transition and its antecedents}
Demkowicz-Dobrza\'nski, Ko\l ody\'nski and Gu\c t\u a
\cite[Eq.~(4) and the dephasing example]{DKG67} use classical channel
simulation to derive metrological bounds, including the noise factor
$\sqrt{1-\lambda^2}/\lambda$ for dephasing (their visibility notation
is different). Thus neither the noisy Heisenberg-to-square-root
transition nor common convex simulation is claimed as a new principle.
Our Lemma~\ref{lem:noisyupper73} instead keeps an exact two-point
programme and finite-$N$ fidelity expression. A proof by finite
entangled blocks gives the reverse modulus, without turning a local
estimation bound into a global distance assertion.

Theorem~\ref{thm:noisycoding73} counts legal reusable descriptions at
all public visibilities, including sequences approaching zero or one.
Its constants are uniform in noise, horizon and small error. Circle
charts give an exact rational upper code, and metric packing allows
centres outside the equatorial family. The conditional-state corollary
combines this noise-dependent direction with the retained rank-state
geometry. These are the asserted covering and realization consequences;
no independent exhaustive priority clearance is implied by identifying
their ingredients. The new pair-dependent programme is not the
family-uniform ambient seizer required by
Theorem~\ref{thm:seizablecover72}. The programme methods of
\cite{DW72,WBHK72,WW72} remain direct antecedents, not inferred
novelty of the present proof.

For a general coupled boundary, the remaining question is how a
nonidentifiable readout or a non-product support change modifies these
operational scales. The equal-row collapse remains an exact warning
against counting a basis that cannot be observed. The current
noise-uniform theorem supplies one explicit transverse approach to an
extremal readout, not a classification of all such directions, nor a
large-error consequence of pairwise packing.

\subsection{Joint visibility and boundary contact}
The single-shot white-noise example of Sedl\'ak--Ziman \cite{SZ74}
already permits different visibilities. The single-shot projective
measurement theorem of \cite{PPKK74} and the multiple-shot theorem
used in Theorem~\ref{thm:readoutentropy72} have different hypotheses and
conclusions. In particular, neither the one-use formula nor the
projective multiple-use formula is an exact noisy-ball formula.
Theorem~\ref{thm:ballmetric74} is instead a two-sided comparison with
absolute constants, proved by a radial support-cutoff argument and
entangled angular blocks. Its covering consequence counts both target
coordinates, keeps unrestricted legal centres in the converse, and
identifies the additional logarithmic factor for the disk.
The contact criterion depends on explicit regularity and boundary-depth
hypotheses; it does not classify every instrument tangent cone.
